\documentclass[journal]{IEEEtran}
\usepackage{cite}
\usepackage{amsmath,amssymb}
\usepackage{url}
\usepackage[T1]{fontenc}

\title{Spatial Accessibility Assessment and Optimal Siting of Social Infrastructure:\\ A Literature Review}

\author{Daulet~Dandykulov~and~Dias~Bolatov\\
        BSc Big Data Analysis (6B06103), Astana IT University\\
        \textit{Supervisor: Aidana Zhalgas}}

\begin{document}

\maketitle

\section{Measuring Spatial Accessibility}
\label{sec:measure}

The simplest accessibility measures partition space administratively and compute a supply-to-demand ratio within each partition, for example places per thousand children in a district. Such container measures are attractive because they use only data that administrations already collect, but they assume that residents use facilities inside their own administrative unit and no others, which is rarely true and is sensitive to the arbitrary geometry of the partition. Buffer or isochrone measures improve on this by defining accessibility relative to the household: a location is served if a facility lies within a fixed distance or travel time, which is the logic embedded in most national planning standards. Buffer measures are transparent and map directly onto regulatory compliance, but they are binary and capacity-blind, treating a location near one overcrowded facility the same as a location near three half-empty ones. The floating catchment area family of measures was developed to combine the household-centred logic of buffers with the supply-and-demand logic of container measures, and it is the dominant framework in the literature reviewed here. The method proceeds in two passes: a catchment is drawn around each supply location and the demand falling inside it is summed to compute a supply-to-demand ratio for that facility, and then a catchment is drawn around each demand location and the ratios of every facility falling inside it are summed, giving an accessibility score per demand location that reflects both how much capacity is reachable and how many other people compete for it. Two parameter choices dominate the behaviour of any such model and are not determined by the data itself: the shape of the distance-decay function, and the size of the catchment. Neither choice is neutral. A preschool's attendance falls off far more sharply with distance than a specialised clinic's does, and impedance itself is not constant through the year in a continental-winter city, where an average pedestrian needs measurably longer to cover the same distance on snow or ice than on bare pavement. Where a national planning standard defines the catchment threshold, that threshold is not a free parameter to be tuned but a regulatory quantity to be respected, which is precisely the design constraint the two technical studies discussed below work within.

A systematic methodological review searching EMBASE, CINAHL, Web of Science and Scopus identified 43 methodological studies and 346 empirical applications of gravity-model-type spatial access measures, and found that almost every such measure in use today descends from the two-step floating catchment area (2SFCA) method or the related kernel density method, with extensions covering distance decay within the catchment, variable catchment size, provider competition, subgroup-specific access, multiple transport modes and time-dependent access. The most frequently used approach in empirical practice, however, remains the simple, original 2SFCA formulation, even as the methodological literature itself has moved toward greater complexity, and most applications still use the resulting measure descriptively rather than as an explanatory variable for outcomes \cite{stacherl2023}. This review is the starting point for the accessibility measure adopted in our own project, since it shows which extensions -- multiple modes, provider competition, decay shape -- are established, defensible choices rather than untested novelties, and it warns that a measure chosen for methodological sophistication alone may not reflect how the measure is actually used and interpreted in practice. A complementary review of place-based accessibility measures examined two of the most recent extensions in particular depth: multimodal accessibility, computed separately for walking, public transport and driving, and temporal change in accessibility, computed at finer temporal granularity than a single annual figure. It proposed two research agendas going forward, improving measurement accuracy through such dynamic variables, and linking that improved accuracy more directly to policy outcomes, since earlier reviews had concentrated mainly on empirical findings, distance-decay functional forms and fixed threshold travel times \cite{park2021}. It shapes our understanding by identifying time-varying impedance, which matters for a city such as Astana with a long, distinct continental winter, as a current research frontier rather than a solved methodological problem; a seasonal, winter-adjusted variant of the pedestrian travel-time matrix is accordingly part of our own project's design.

Two technical studies apply exactly these extensions and are directly relevant to how our own accessibility measure should be built. Mao, Li and Yan address the problem that standard accessibility models do not adequately capture different transport modes, facility attractiveness or user choice behaviour, and propose CM2SFCA, a comprehensive multimodal 2SFCA formulation for measuring accessibility to parks in Wuhan, China, that combines facility capacity, population demand, distance decay, variable catchment areas and multiple transport modes within a single model. The resulting model gives a more realistic accessibility picture than computing cycling, driving and public-transport accessibility separately, and it successfully identifies low-accessibility areas across a range of population densities; its contribution is showing that all of these factors can be integrated into one 2SFCA formulation rather than treated as separate, disconnected calculations. Its main limitation is that the method is considerably more complex than the standard approach and needs extensive transport and spatial data, and it is demonstrated on parks rather than on preschools or schools, which is the facility type our own project addresses \cite{mao2023wuhan}. Yhee, Kim and Kang address a related but distinct problem: that standard social-infrastructure planning based on population counts per area does not capture how convenient it actually is for residents to reach a facility. Their method combines a GIS-based accessibility index with real travel times, by driving and walking, obtained from a navigation API rather than straight-line distance, and it identified specific areas in South Korea with inadequate access to libraries, parks, schools, kindergartens and sports facilities. Its contribution to knowledge is substituting routed travel time for straight-line distance across several facility types at once, which advances the field beyond the Euclidean-distance assumption still common in container and buffer measures; its limitations are that it does not account for facility capacity, excludes residents outside the study boundary, and selects the nearest facility purely by distance rather than by shortest travel time, which can misidentify the facility a resident would actually use. It is directly relevant to our own use of real, routed pedestrian travel times, rather than straight-line distance, when computing accessibility across Astana's road and pedestrian network \cite{kim2021gis}.

\section{Accessibility, Equity and Applied Case Studies}
\label{sec:equity}

An accessibility surface is, on its own, only a map. To make claims about fairness it must be reduced to statistics that describe how unequally accessibility is distributed across people, not merely across space: a deficit in an empty industrial zone is not an equity problem in the way that the same deficit is in a dense residential district. The standard instruments for this reduction are borrowed from income-inequality analysis: the Lorenz curve plots the cumulative share of accessibility against the cumulative share of population, ordered from least to most advantaged, and the Gini coefficient summarises the area between that curve and the line of perfect equality into a single number. A scalar index, however, says nothing about \emph{where} inequality is concentrated, and a distribution can appear acceptably equal in aggregate while a specific, identifiable subpopulation is systematically disadvantaged. Evaluating equity properly therefore requires attention to how accessibility is distributed over particular areas and demographic groups, not only to its aggregate spread across the city as a whole, which is exactly the level of detail the studies reviewed in this section were designed to provide. Several studies in this review apply accessibility measurement specifically to populations for whom the distinction between aggregate provision and lived access is most consequential, which helps frame who the equity analysis in our own project should focus on: children of preschool and school age, and, by analogy, the elderly and rural populations examined in the adjacent literature.

A systematic review of sixteen studies published since 2000 on older adults' accessibility to destinations by public transport found that accessibility metrics, destination types, transport modes and the conceptualisation of older people vary so widely across the literature, with no single method dominating, that cross-study conclusions are difficult to draw. It called for research that engages with critical gerontology and age-friendly city frameworks, involving older people directly in choosing which destinations and thresholds matter to them, for example through experience-based co-design rather than analyst-imposed assumptions \cite{ravensbergen2022}. It shapes our understanding by showing that results depend heavily on which destinations and thresholds are chosen, so our own project must justify its thresholds explicitly, anchored on the statutory service radii defined in the Kazakh planning standard rather than on values borrowed from an unrelated study, and must treat subgroup-specific accessibility, such as accessibility disaggregated by age cohort, as a first-class analytical concern rather than an afterthought reported only for the population as a whole. A parallel systematic review of geospatial approaches to educational inequity in rural communities screened 1247 records down to the 16 that met its inclusion criteria, and found that distance modelling, network travel-time analysis, cost-distance modelling, spatial autocorrelation, kernel density estimation, spatial regression and multi-criteria opportunity mapping are the methodologies most frequently used in this field. These techniques revealed patterns of disadvantage, in school accessibility, teacher and infrastructure shortages, and clustering of low opportunity, that non-spatial approaches missed entirely \cite{rwandashi2026}. Because its focus is specifically rural communities, it confirms that the urban, rapidly growing setting of a capital such as Astana is not directly covered by this literature, a geographic gap that our own project is designed to address.

Two technical case studies illustrate exactly this kind of disadvantage empirically and are directly relevant to the equity dimension of our own project. A GIS-based spatial analysis of upper secondary education access in Italy's Campania region addresses the problem that rural and remote residents have markedly poorer geographic access to schools than urban residents, with social and infrastructural barriers compounding the effect. Using census-tract-level data and transport-accessibility metrics, it found that rural students face substantially longer commutes to school, and that low household income and poor internet coverage further exacerbate the disadvantage rather than substitute for it. Its contribution to knowledge is empirical evidence that physical distance is a barrier to education access in its own right, and not merely a proxy for economic disadvantage that would disappear once income and connectivity are controlled for; its limitation is that the study covers only one region and one school level, so its findings should not be generalised directly to a different regulatory or geographic context \cite{bruno2022italy}. It supports our project's treatment of physical accessibility as a distinct equity dimension, separate from and additional to income-based measures of disadvantage. In Chengdu, a study addresses the problem of whether existing school locations in a fast-growing mega-city actually align with where children live, rather than with where they were historically built. It combined a Random Forest model of the spatial distribution of school-age children with kernel density estimation, an imbalance (Lorenz-type) index, a multimodal MM2SFCA accessibility measure and Moran's~$I$ spatial autocorrelation analysis, and found a decentralised, multi-core pattern of child population, a significant imbalance in school placement ($S = 0.257$), and specific ``low accessibility, high density'' pockets concentrated in the city's second ring, outside the historic core but not yet served by new construction. Its contribution to advancing knowledge is combining demographic machine-learning modelling with multimodal accessibility analysis to assess social equity directly and quantitatively, rather than relying on population totals alone; its limitation is that it does not account for school quality or reputation, both of which shape real demand and household choice independently of geographic distance. It is a direct precedent for our own plan to disaggregate WorldPop population data by age cohort before computing accessibility, applied to a city, Astana, that is growing at a comparably fast rate \cite{jiang2024chengdu}.

\section{Facility Location and Optimization}
\label{sec:location}

Where accessibility measurement is descriptive, facility location theory is prescriptive: given a demand surface and a budget of new facilities, where should they be placed to do the most good? This is the second half of what our project needs to deliver, and it draws on a mature body of operations-research literature built around two competing objectives. The $p$-median model selects a fixed number of facility locations so as to minimise the total demand-weighted travel distance across the whole population, and is the natural formulation when the planning objective is average convenience; it will, however, happily improve the situation of already well-served residents if doing so lowers the city-wide average, and is therefore indifferent to whether an already-underserved area is helped at all. The maximal covering location problem (MCLP) instead maximises the amount of demand falling within a fixed coverage radius for a given budget of facilities, and is indifferent to improvements inside the already-covered set, rewarding only the conversion of previously uncovered demand into covered demand. Where a planning standard defines a statutory service radius, as it does in Kazakhstan, MCLP is the formulation that matches the regulatory objective directly, and a capacitated version of the model, which additionally respects that a facility cannot serve unlimited demand, is required whenever capacity data are available, since school and preschool places are the binding constraint in practice. An overview written for the 25th anniversary of \emph{Transportation Research Part E} traces the long relationship between facility location and logistics and transportation, revisiting conventional location problems -- among them the classical $p$-median and covering formulations -- alongside newer areas such as reverse and green logistics, humanitarian supply chains and multimodal transportation, and identifies uncertainty modelling, and the shift from Industry 4.0 toward data-driven Industry 5.0 decision making, as the field's main current challenges \cite{saldanha2022}. It frames location optimisation as an inherently data-driven, uncertainty-aware task rather than a one-shot deterministic exercise, which supports our own plan to report a sensitivity analysis of demand and parameter assumptions rather than a single point estimate presented as though it were certain. A broader survey of facility location problems over roughly the past decade categorises the literature by problem variant, mathematical formulation, solution method and exact approach, the methodologies most frequently used in this field, and reports steady progress in mathematical programming, scalable exact methods and hybrid metaheuristics, with uncertainty modelling, multi-objective optimisation and sustainability named as the field's emerging research gaps \cite{dissanayake2024}. Although its setting is logistics rather than social infrastructure, it supports our decision to report a coverage-versus-equity trade-off, for example as a Pareto frontier, rather than collapsing the two competing objectives into a single, arbitrarily weighted score that hides the trade-off from the decision maker.

A study of rural Changyuan, China, addresses the problem that access to schools in rural areas is uneven, with some territories left completely unserved by any school facility, and it shows what applying facility-location theory to social infrastructure specifically can look like from end to end. It measured accessibility with the MH3SFCA method, then optimised school locations with a capacitated maximal covering model (MCLP-CC), and compared the resulting situation before and after optimisation using the same accessibility measure throughout, so that the improvement is directly attributable to the siting decision rather than to a change in methodology. After optimisation, the unserved area fell by 62.55\% for junior high schools, 70\% for elementary schools, and 84.30\% for nursery schools. Its contribution to advancing knowledge is integrating accessibility assessment, optimisation and validation into a single, internally consistent planning framework, showing concretely that a measure-then-optimise-then-validate pipeline can be built and its improvement quantified end to end, rather than left as separate, disconnected analytical steps; its limitation is that it is set in a rural area and focuses specifically on schools, so applying the same framework to other facility types or to a dense urban setting needs separate verification before the results can be assumed to transfer \cite{han2023education}. Among all twenty sources reviewed here, this is the closest precedent for our own project's evaluation-optimisation-validation structure, and specifically for backtesting an optimiser's recommendations against the realised siting decisions of Astana's national Comfortable School construction programme.

\section{Machine Learning, Explainability and Site Selection}
\label{sec:ml}

Gradient-boosted decision tree ensembles, in implementations such as XGBoost, LightGBM and CatBoost, are now routinely used for tabular and spatial urban data because they handle heterogeneous features and non-linear feature interactions without manual specification, tolerate missing values gracefully, and train quickly enough to permit thorough cross-validation even on moderately sized datasets. Because the resulting models are not directly interpretable, their adoption in planning contexts has been accompanied, in the literature reviewed here, by post-hoc explanation methods, and this section covers both that methodological literature and the applications closest to our own siting model. A systematic review of 581 papers on spatially-explicit geospatial artificial intelligence (GeoAI) in urban geography, organised into urban dynamics, social differentiation of urban areas, and social sensing, found that deep-learning applications in this space are proliferating rapidly, but that spatially-explicit GeoAI models specifically are still at an early stage of development. It identified three specific challenges of existing models as open research gaps, and called for future work to develop multi-scale, explainable, spatially-explicit GeoAI rather than treating spatial structure as an afterthought bolted onto a generic model \cite{liu2022}. It supports our decision to treat explainability as a design requirement rather than an optional extra for any machine-learning component of the project. A broader review of machine learning across geospatial applications, spanning precision agriculture, environmental monitoring, disaster management and urban planning, covers a wider scope than the urban-geography-specific review above; it presents GeoAI as the systematic combination of AI and ML models with geospatial data, and emphasises its links to GIS and to large-scale ``big geodata'' infrastructure. Its methodological survey and stated future prospects confirm that machine learning on geospatial data is now a mainstream, fast-growing area that explicitly includes urban planning as an application domain, though the review's breadth means it offers little guidance specific to accessibility measurement or facility siting \cite{boutayeb2025}.

On the specific problem of interpretability, that machine-learning models are often highly accurate but function as a ``black box,'' making it difficult to interpret the spatial effects they have implicitly learned, a study combining XGBoost gradient boosting with the SHAP explanation method, tested both through simulation against a Spatial Lag Model and a Multi-scale Geographically Weighted Regression model and through an empirical case study of Chicago ride-hailing data, found that SHAP values assigned to spatial coordinates accurately reflect spatial effects and align closely with estimates from classical spatial econometrics. Its contribution to advancing knowledge is combining the predictive accuracy of machine learning with the interpretive transparency traditionally associated only with spatial statistical models, effectively bridging two previously separate methodological traditions; its main limitation is dependence on hyperparameter choices and sensitivity to spatial multicollinearity among predictors, which can distort attribution if not checked \cite{li2022extracting}. It gives the direct methodological justification for interpreting our own gradient-boosted siting model with SHAP, rather than fitting a separate, purpose-built spatial regression model that would require additional specification choices of its own. A related benchmark addresses a narrower but connected problem, that location encoders convert raw coordinates into learned embeddings whose ``black-box'' nature disperses geographic information across all embedding dimensions, making the spatial effects those embeddings capture difficult to interpret directly. It evaluated eleven location encoders from the TorchSpatial framework using GeoShapley on synthetic data with known ground truth, across three spatial scales, and found that the primary coefficient is recovered very well ($r > 0.98$) while recovery of the secondary coefficient depends on scale, with methods such as Sphere2Vec-sphereM+ performing best at the global level and contrastive training often degrading rather than improving performance. Its contribution is being the first evaluation of GeoShapley's ability to extract interpretable spatial effects specifically from task-agnostic neural coordinate encoders, rather than from simpler statistical or tree-based models; its limitation is that the evaluation uses smooth synthetic data with a fixed embedding dimension tied to one particular MLP architecture, so its conclusions are best treated as a starting point for selecting encoder architectures rather than as a general, architecture-independent result \cite{li2024geoshapley}. It is a candidate refinement for how we interpret our own siting model's SHAP attributions, since it would let us separate a genuine locational effect from covariates that merely happen to correlate with geography without themselves being spatial.

Two studies apply these explainable machine-learning methods directly to school siting, and are the closest published precedents for our own project's central analytical step. In Tianjin, addressing the problem of how schools and kindergartens can be equitably distributed, keeping every household within a 15-minute walk, as the city continues to expand, a 500$\times$500\,m grid analysis of points of interest, population density and transport data compared classical statistical models against Random Forest and XGBoost, interpreted afterward with SHAP to identify which factors drove the predictions. The ensemble models exceeded 85\% accuracy, reaching 92.41\% in one configuration, with child-population density, transport accessibility and proximity to other schools emerging as the leading factors, and specific priority development zones were identified for new construction. This constitutes an automated, data-driven forecasting tool for urban planning built on large administrative and geospatial datasets; its principal limitation, in our assessment, is that it relies on the historical distribution of existing schools and lacks fine-grained local indicators, so the model in fact describes where schools were placed historically rather than proving where they should be placed going forward. It is the closest methodological precedent for our own siting model, and the specific study against which our project's revealed-preference interpretation, reading such a model as evidence about past planning decisions rather than as normative guidance, is explicitly positioned \cite{sun2025tianjin}. In \.{I}zmir, addressing the problem of optimally selecting sites for new school construction so as to ensure an equitable distribution of educational services and support sustainable urban development, a hybrid GIS and machine-learning framework evaluated 12 criteria across 39 candidate sites using the Analytic Hierarchy Process, equal weighting, XGBoost-derived weights and TOPSIS, alongside a custom machine-learning ranking method (ML\_Rank) developed for the study. It identified site AS4 as the most suitable location overall, and found that ML\_Rank agreed closely with equal weighting ($\rho = 0.9745$) and with XGBoost-derived weights ($\rho = 0.8813$), while the expert-elicited AHP weights showed the weakest agreement with the other three methods. This demonstrates that machine-learning-derived criterion weights can be more objective and internally consistent than traditional expert-based weighting once the criteria set grows large and complex; its limitation is that the resulting ranking is confined to two districts and one fixed set of 12 criteria, so its transferability to a different regulatory context is untested \cite{basegmez2025izmir}. It is a recent, concrete example of combining GIS, multi-criteria decision analysis and machine learning for ranking a discrete set of candidate sites, directly comparable to the candidate-site ranking step of our own project's optimisation stage.

\section{Validating Spatial Models}
\label{sec:validation}

A recurring methodological concern across this entire literature is whether standard model-evaluation procedures can be trusted when the underlying data are spatially dependent, which matters directly for how our own siting model should be validated before its results are reported as reliable. Nearby locations resemble one another, a property known in geostatistics as spatial autocorrelation and often summarised informally as Tobler's first law of geography, that everything is related to everything else, but near things are more related than distant things. As a consequence, a randomly selected test set commonly contains observations that are near-duplicates of observations already seen during training, and a model evaluated on such a test set can appear to perform far better than it actually would on a genuinely new, geographically distinct area, since the model is effectively being re-tested on data it has already, in a spatial sense, seen. This has direct bearing on the machine-learning siting studies reviewed in the previous section: a reported accuracy above 85\% for a facility-suitability model is difficult to interpret without knowing how the validation folds were constructed, since the target variable in such models, the presence of a facility, is by construction strongly spatially clustered. One study addresses the specific problem of whether spatial cross-validation, as commonly practised, accurately assesses map accuracy; through a methodological and statistical argument rather than a new empirical data-collection exercise, it shows that both standard and spatial cross-validation can give biased estimates of map accuracy, and that only probability sampling combined with design-based statistical inference is demonstrably reliable. Its contribution to the field is debunking the common misconception that spatial cross-validation is, by itself, sufficient to obtain an unbiased evaluation of map accuracy; its limitation, from our perspective, is that it is written for maps of ecological and environmental variables rather than for urban infrastructure, where a probability sample of accuracy-check locations is simply not available given how facility data are collected \cite{wadoux2021spatial}. It justifies our use of spatially blocked cross-validation for our own siting model as the best practically available option, while honestly requiring us to acknowledge, exactly as this source does, that unbiasedness in the strict statistical sense is still not thereby guaranteed.

A review of machine-learning studies that predict environmental variables while explicitly accounting for spatial autocorrelation identified four families of method that are the methodologies most frequently used in this field: explicit spatial covariates added directly to the model, spatially aware train/test splitting that keeps nearby points from appearing on both sides of a fold, model-level statistical calculations, and independent exploratory analysis conducted separately from the modelling pipeline itself. It found that authors most often handle autocorrelation only during that final, independent exploratory stage, at a point where it can no longer have any effect on the actual predicted values, and it identifies the field's main outstanding research gap as the absence of an overall systematic approach, with the appropriate method judged to depend heavily on the specific study and dataset rather than following a settled convention \cite{jemeljanova2024}. It shapes our understanding by supporting spatially aware data splitting specifically for our own siting model, and by warning explicitly that whichever splitting strategy we choose must be justified for our particular case rather than adopted by default, since the studies it reviews are drawn from environmental science rather than from urban planning applications. A more explicitly statistical treatment of the same underlying problem reviews regularised estimation methods for geostatistical and spatial econometric models alongside cross-validation techniques specifically designed for spatiotemporal data, and sets out three key assumptions that any valid partitioning of the data must satisfy across time-series, spatial and spatiotemporal structures, together with a discussion of available software implementations for practitioners \cite{otto2024}. It gives explicit, citable criteria for designing the spatial folds used to validate our own siting model, although it is written primarily for geostatistical and econometric models rather than for gradient-boosted trees, which somewhat limits how directly its specific guidance transfers to our own method and means its criteria must be adapted rather than applied verbatim.

\bibliographystyle{IEEEtran}
\bibliography{refs}

\end{document}
